\section{Hyperbolicity and outer automorphisms}
\label{hyp:section}

We now prove that the common manifold group \(\Pi\) has finite outer
automorphism group.  The geometric step is to show that the reflection
extension \(\mathcal E=W\rtimes H\) is word-hyperbolic.  We use the
support construction of \cite[Lemma~3.3 and its proof]{PD} to obtain a
cubical model for \(\mathcal E\), and then bound the size of isometric
square grids in its edge graph.  Poincar\'e duality will exclude the
virtually cyclic splittings that an infinite outer automorphism group
would force.

Write \(\widehat{\mathcal V}=\widehat{\mathcal L}^{(0)}\) for the
generating set of \(W\).  Recall that \(\widehat{\mathcal L}\) is a
locally finite flag triangulation of a PL three-manifold, and that its
graph has no induced four-cycle.  In particular, every clique has at
most four vertices.  The deck group \(H\) acts freely on
\(\widehat{\mathcal L}\) with finite quotient.

\subsection{The cubical model and its edge graph}

The ordinary Davis complex of \(W\) has infinitely many generator
types.  The following construction, recalled from
\cite[Lemma~3.3]{PD}, arranges Davis complexes of finitely generated
special subgroups over a tree on which \(H\) acts.  We record the
resulting edge graph because the hyperbolicity argument will use its
labels directly.

\begin{lemma}[Support construction]
\label{hyp:cubical-model}
There are a locally finite tree \(T\) with a free cocompact \(H\)-action,
an \(H\)-equivariant family of finite subtrees
\((T_u)_{u\in\widehat{\mathcal V}}\), and a locally finite CAT(0) cube
complex \(X\) with a proper cocompact \(\mathcal E\)-action, with the
following properties:
\begin{enumerate}
\item If \(u,u'\) are adjacent in \(\widehat{\mathcal L}\), then
      \(T_u\) and \(T_{u'}\) meet at a vertex.  Moreover,
      \[
      b_T:=\sup_{u\in\widehat{\mathcal V}}|T_u^{(0)}|<\infty.
      \]
\item The vertices of \(X\) are the pairs \((a,t)\in W\times T^{(0)}\).
      Its edges are exactly
      \begin{equation}
      \begin{aligned}
      (a,t)&\longleftrightarrow(au,t)
        &&\text{if }u\in\widehat{\mathcal V}\text{ and }t\in T_u^{(0)},\\
      (a,t)&\longleftrightarrow(a,t')
        &&\text{if }[t,t']\text{ is an edge of }T.
      \end{aligned}
      \label{hyp:edge-graph}
      \end{equation}
\item On these vertices the action is
      \[
      (b,h)(a,t)=(b\,h(a),ht)
      \qquad (b\in W,\ h\in H).
      \]
      The cubes of \(X\) have dimension at most five.
\end{enumerate}
\end{lemma}

\begin{proof}
The group \(H\) is finitely generated: its free action on the connected
locally finite graph \(\widehat{\mathcal L}^{(1)}\) has finite quotient,
so lifting a spanning tree of that quotient leaves only finitely many
edge identifications needed to generate the deck group.  Choose the
Cayley tree \(T\) of a finite free basis of \(H\).  It is locally finite,
and \(H\) acts freely and cocompactly without inversions.

Choose points \(a_u\in T^{(0)}\), equivariantly in
\(u\in\widehat{\mathcal V}\); one chooses them on representatives of
the finitely many \(H\)-orbits and then translates.  Freeness of the
deck action makes this assignment well defined.  There are also only
finitely many orbits of edges of \(\widehat{\mathcal L}\).  Thus
\(d_T(a_u,a_{u'})\), for adjacent \(u,u'\), has a common finite bound.
Choose an integer \(r\geq 0\) at least this bound and let \(T_u\) be the
closed combinatorial ball of radius \(r\) about \(a_u\).  These balls
are subtrees.  Adjacent vertices have supports meeting at a vertex,
and the number of vertices in each support is bounded uniformly by a
finite number \(b_T\).  In particular \(b_T\geq 1\).

For a vertex \(t\) and an edge \(e\) of \(T\), define
\[
\mathcal U_t=\{u:t\in T_u\},\qquad
\mathcal U_e=\{u:e\subset T_u\}.
\]
Both kinds of set have uniformly bounded finite size.  Indeed, a
support containing \(t\) has its center in the radius-\(r\) ball about
\(t\); that ball contains uniformly finitely many orbit points, and
there are only finitely many generator orbits.  The assertion for
\(\mathcal U_e\) follows from the assertion for either endpoint.
Since each \(T_u\) is a subtree, for \(e=[t,t']\) we have
\begin{equation}
\mathcal U_e=\mathcal U_t\cap\mathcal U_{t'}.
\label{hyp:endpoint-supports}
\end{equation}
Let \(W_t\) and \(W_e\) be the special subgroups of \(W\) generated by
\(\mathcal U_t\) and \(\mathcal U_e\), respectively.  The edge maps
\(W_e\longrightarrow W_t,W_{t'}\) are special subgroup inclusions.

These groups form a graph of groups over \(T\) whose fundamental group
is \(W\).  To see the presentation explicitly, the copies of a
generator \(u\) in the vertex groups are identified along precisely
the connected subtree \(T_u\), by \eqref{hyp:endpoint-supports}.
Every commutation relation of \(W\) occurs in a vertex group, since
the supports of its two adjacent generators meet at a vertex.  Each
vertex group is defined by the induced subgraph on \(\mathcal U_t\),
so it introduces only relations of \(W\).  The resulting presentation
is therefore the right-angled presentation of \(W\).  Special
subgroups inject, so the Bass--Serre tree of this graph of groups has
vertices \((aW_t,t)\) and edges \((aW_e,e)\), with the evident endpoint
maps.  This Bass--Serre tree, which can have infinite valence, is the
tree over which the next spaces are assembled.

For a finite set \(J\subset\widehat{\mathcal V}\), let \(\Sigma_J\)
denote the Davis cube complex of \(W_J\).  Its vertices are the elements
of \(W_J\).  For each clique \(C\subset J\) and each coset \(kW_C\)
in \(W_J\), it has a cube whose vertex set is \(kW_C\); its dimension
is \(|C|\), since \(W_C\cong(\mathbb Z/2)^{|C|}\).  For \(J=\varnothing\)
this is a single vertex.  These complexes are CAT(0), and inclusions
of special subgroups induce closed convex cubical inclusions of their
Davis complexes; see \cite[Chapters~7 and~12]{Davis2008}.  Set
\(\Sigma_t=\Sigma_{\mathcal U_t}\) and
\(\Sigma_e=\Sigma_{\mathcal U_e}\).

Following the support construction of \cite[Lemma~3.3]{PD}, use the
vertex spaces \(W\times_{W_t}\Sigma_t\) and the edge cylinders
\[
W\times_{W_e}\bigl(\Sigma_e\times[0,1]\bigr).
\]
Here \(W\times_K Z\) identifies \((ak,z)\) with \((a,kz)\) for
\(k\in K\).  Choose an ordering \(e=[t,t']\) of the endpoints of each
edge of \(T\), and glue the cylinder at \(0\) and \(1\) into the
corresponding vertex spaces by the special-subcomplex inclusions.
Give \([0,1]\) its single-edge cubulation.  The resulting cubical
complex \(X\) is a tree of spaces over the Bass--Serre tree just
described.

We can now identify its graph.  A vertex of \(W\times_{W_t}\Sigma_t\)
is represented by \([a,k]\) with \(k\in W_t\), and
\([a,k]\mapsto ak\) is a bijection from these vertices to \(W\).
An edge in \(\Sigma_t\) joins \(k\) to \(ku\) for
\(u\in\mathcal U_t\), so the vertex-space edges are exactly the first
line of \eqref{hyp:edge-graph}.  The same identification gives one
vertex labeled by each element of \(W\) in
\(W\times_{W_e}\Sigma_e\).  Each endpoint inclusion preserves this
label.  Thus its cylinder has an edge from \((a,t)\) to \((a,t')\)
for every \(a\in W\).  The cylinder has no vertices over the interior
of \([0,1]\); its remaining edges lie in its endpoint spaces and are
already vertex-space edges.  This proves that
\eqref{hyp:edge-graph} lists every edge, including when
\(\mathcal U_e=\varnothing\).

For completeness, the local and metric properties in the cited
construction are compatible with this graph description.  For a fixed
edge \(e\) incident to \(t\), the translates of \(\Sigma_e\) in a
vertex copy of \(\Sigma_t\) have as their vertex sets the cosets of
\(W_e\) in \(W_t\).  Those cosets partition the Cayley vertices.
Hence a finite subcomplex meets only finitely many such translates,
even if the Bass--Serre vertex has infinite valence.  There are only
finitely many edges of \(T\) incident to \(t\), and \(\Sigma_t\) is
locally finite.  It follows that \(X\) is locally finite.  A vertex
cube has dimension at most four and an edge cylinder at most five,
by the clique bound.  The unit-cube length metric is consequently
proper and complete.  Convex CAT(0) gluing first over finite subtrees
of the Bass--Serre tree, and then over their union, proves that \(X\)
is CAT(0): the finite-subtree gluings are CAT(0) and convex in larger
ones, and completeness follows from the local finiteness just checked.

Left multiplication defines the \(W\)-action.  Equivariance of the
supports defines the \(H\)-action on the vertex and edge spaces, and
on vertices the combined action is
\((b,h)(a,t)=(b\,h(a),ht)\).  This action is free on vertices:
if it fixes \((a,t)\), then \(h\) fixes \(t\), whence \(h=1\), and
then \(ba=a\) gives \(b=1\).  Local finiteness makes the cellular
action proper.  There are finitely many vertex orbits, because \(W\)
acts transitively on the first coordinate and \(T/H\) is finite.
Local finiteness then gives finitely many cube orbits, so the action is
cocompact.
\end{proof}

\subsection{A uniform bound on square grids}

Give \(X^{(1)}\) its edge-path metric \(d_1\).  Label an edge in the
first line of \eqref{hyp:edge-graph} by its generator \(u\), and label
an edge in the second line by the unoriented edge \([t,t']\) of \(T\).
We call these Coxeter labels and tree labels, respectively.  For an
integer \(n\geq0\), a square grid of side length \(n\) means the graph on
\(\{0,\ldots,n\}^2\) with its \(\ell^1\) metric.

\begin{lemma}
\label{hyp:bounded-grids}
Every isometric embedding of a square grid of side length \(n\) into
\(X^{(1)}\) satisfies
\[
n\leq 16b_T-1.
\]
\end{lemma}

\begin{proof}
First consider any embedded four-cycle in \(X^{(1)}\).  Projection to
\(T\) shows that a cycle containing tree edges has either two or four
of them.  Four tree edges would give an embedded cycle in a tree,
because the \(W\)-coordinate would be constant.  With two tree edges,
the projected closed walk traverses the same tree edge twice.  They
cannot be consecutive in the four-cycle, since two such consecutive
steps would repeat a vertex.  They are therefore opposite.  The two
remaining edges are also opposite and have the same Coxeter label,
as their product in \(W\) is the identity.

If all four edges have Coxeter labels, their label word is trivial in
\(W\).  The abelianization of \(W\) is the direct sum of copies of
\(\mathbb Z/2\) indexed by \(\widehat{\mathcal V}\), so each label
occurs an even number of times.  Embeddedness forbids consecutive
equal labels, including at the closing vertex.  Thus two distinct
labels alternate.  They commute: if they were nonadjacent, the
special subgroup on them would be the infinite dihedral group, in
which their alternating word of length four is nontrivial.  We have
proved that opposite edges of every embedded four-cycle have equal
labels, adjacent edges cannot both have tree labels, and adjacent
Coxeter labels are distinct and commute.

Let an isometric grid of side length \(n\) be given.  Each of its
elementary four-cycles is embedded.  Equality of labels on opposite
edges shows that the label at a fixed horizontal step is independent
of the vertical coordinate, and the label at a fixed vertical step is
independent of the horizontal coordinate.  If each direction had a
tree label, the elementary square where those steps meet would have
adjacent tree edges.  After interchanging the directions if necessary,
all vertical labels are therefore Coxeter labels.

Suppose \(n\geq16\).  The vertical labels include two that do not
commute.  Indeed, if they all commuted, their distinct labels would
form a clique of size at most four.  They would generate a group of
order at most \(2^4=16\).  Each column has constant tree coordinate,
so its \(n+1\) distinct vertices would have at most 16 possible
\(W\)-coordinates, a contradiction.

Every horizontal Coxeter label is distinct from and commutes with
every vertical label, by the elementary-square calculation.  If two
horizontal Coxeter labels did not commute, they and the two
noncommuting vertical labels would be four distinct vertices whose
cross pairs are adjacent and whose two same-direction pairs are
nonadjacent.  They would form an induced four-cycle in
\(\widehat{\mathcal L}\).  Hence the horizontal Coxeter labels form a
clique and generate a group of order at most 16.  Along any one row,
tree edges leave the \(W\)-coordinate unchanged, so all
\(W\)-coordinates lie in a single left coset of this finite group.

Fix one vertical label \(u\).  Each column has constant tree
coordinate, and the \(u\)-edge occurs in every column at the same
vertical step.  The edge description \eqref{hyp:edge-graph} therefore
places every tree coordinate of the row in \(T_u^{(0)}\).  There are
at most \(b_T\) such coordinates and at most 16 \(W\)-coordinates.
The row has \(n+1\) distinct vertices, so
\[
n+1\leq16b_T.
\]
This proves the claimed bound when \(n\geq16\).  If \(n\leq15\), the
same bound follows from \(b_T\geq1\).
\end{proof}

\subsection{From grids to hyperbolicity}

We next explain why this grid bound implies thin triangles.  We use
the standard hyperplane description of the median graph of a CAT(0)
cube complex, as in
\cite[Theorem~6.1, Lemmas~6.3--6.4, and the paragraph following
Lemma~6.4]{Chepoi}.
For vertices \(x,z\), let \(J_x(z)\) be the finite set of hyperplanes
separating them.  If \(d_1\) denotes the graph metric, then
\begin{equation}
d_1(z,z')=|J_x(z)\mathbin{\triangle}J_x(z')|.
\label{hyp:hyperplane-distance}
\end{equation}
In particular, \(z\) belongs to the interval
\[
I(x,y):=\{z:d_1(x,z)+d_1(z,y)=d_1(x,y)\}
\]
exactly when \(J_x(z)\subset J_x(y)\).  Every triple of vertices has
a median, the unique vertex in all three pairwise intervals.  Its
side of each hyperplane is the side containing at least two of the
three vertices.  Equivalently, for any vertices \(a,b,c\),
\begin{equation}
J_x(\operatorname{med}(a,b,c))
 =\bigl(J_x(a)\cap J_x(b)\bigr)
  \cup\bigl(J_x(b)\cap J_x(c)\bigr)
  \cup\bigl(J_x(c)\cap J_x(a)\bigr).
\label{hyp:median-sides}
\end{equation}

\begin{lemma}
\label{hyp:grids-to-triangles}
Let \(Y\) be a CAT(0) cube complex.  Suppose every isometrically
embedded square grid in its edge graph has side length at most
\(K\), where \(K\) is a nonnegative integer.  Then any two vertices
of \(I(x,y)\) at the same distance from \(x\) are at distance at most
\(2K\), uniformly in \(x,y\).  Every geodesic triangle with vertex
corners in the edge graph is \(4K\)-thin at its vertices.
\end{lemma}

\begin{proof}
Use the notation and hyperplane facts above for \(Y\).  Fix
\(v,v'\in I(x,y)\) with \(d_1(x,v)=d_1(x,v')\).  Then
\(J_x(v),J_x(v')\subset J_x(y)\) and these two finite sets have equal
cardinality.  Put
\[
B=J_x(v)\cap J_x(v'),\qquad
n=|J_x(v)\setminus J_x(v')|
 =|J_x(v')\setminus J_x(v)|.
\]
For \(m=\operatorname{med}(x,v,v')\), Equation
\eqref{hyp:median-sides} gives \(J_x(m)=B\).
Choose edge geodesics
\[
v_0=m,v_1,\ldots,v_n=v,\qquad
v'_0=m,v'_1,\ldots,v'_n=v'.
\]
Their lengths are \(n\) by \eqref{hyp:hyperplane-distance}.
Along each geodesic, the separating sets add exactly one hyperplane
at each step: its endpoint set contains its initial set, and removing
a hyperplane or adding one outside the endpoint set would require an
extra crossing.  Thus
\[
J_x(v_i)=B\cup A_i,\qquad
J_x(v'_j)=B\cup A'_j,
\]
where \(A_i\) and \(A'_j\) are increasing chains of sets of
cardinalities \(i\) and \(j\), respectively.  All \(A_i\) lie in
\(J_x(v)\setminus J_x(v')\), and all \(A'_j\) lie in the disjoint set
\(J_x(v')\setminus J_x(v)\).

For \(0\leq i,j\leq n\), set
\(z_{ij}=\operatorname{med}(y,v_i,v'_j)\).
Since the two separating sets for \(v_i,v'_j\) are contained in
\(J_x(y)\), the majority formula gives
\[
J_x(z_{ij})=J_x(v_i)\cup J_x(v'_j)=B\cup A_i\cup A'_j.
\]
The two chains add disjoint hyperplanes.  Their nesting and
\eqref{hyp:hyperplane-distance} therefore give
\[
d_1(z_{ij},z_{kl})=|i-k|+|j-l|
\qquad (0\leq i,j,k,l\leq n).
\]
In particular these vertices are distinct and form an isometrically
embedded square grid of side length \(n\).  Hence \(n\leq K\), and
\[
d_1(v,v')=|J_x(v)\mathbin{\triangle}J_x(v')|=2n\leq2K.
\]
This proves the uniform assertion about intervals.

For vertices \(a,b,c\), let \(m\) be their median and choose one edge
geodesic from \(m\) to each corner.  Concatenating the appropriate
two chosen paths gives a geodesic between each pair of corners,
because \(m\) belongs to every pairwise interval.  Each of these
three chosen sides lies in the union of the other two.

Now take arbitrary edge-geodesic sides of the triangle.  A vertex on
the arbitrary side from \(a\) to \(b\) and the vertex at the same
distance from \(a\) on the chosen side both belong to \(I(a,b)\), so
they are at distance at most \(2K\).  The latter vertex lies on the
chosen side from \(a\) to \(c\) or the chosen side from \(b\) to \(c\).
Matching its distance from the corresponding endpoint on that
arbitrary side costs at most another \(2K\).  Thus every vertex on
each arbitrary side is within \(4K\) of the other two sides.  This
vertex formulation is the usual thin-triangle criterion for the
unit-edge graph.
\end{proof}

\begin{proposition}
\label{hyp:word-hyperbolic}
The groups \(\mathcal E\) and \(\Pi\) are word-hyperbolic.
\end{proposition}

\begin{proof}
Lemmas~\ref{hyp:bounded-grids} and \ref{hyp:grids-to-triangles}, with
\(K=16b_T-1\), show that \(X^{(1)}\) is a hyperbolic graph.  Its
unit-edge metric is proper because \(X\) is locally finite.  The
\(\mathcal E\)-action on this graph is proper and cocompact by
Lemma~\ref{hyp:cubical-model}.  The \v{S}varc--Milnor lemma
\cite[Proposition~I.8.19]{BridsonHaefliger1999} gives a quasi-isometry
from a Cayley graph of \(\mathcal E\) to \(X^{(1)}\).
Hyperbolicity is invariant under quasi-isometry
\cite[Theorem~III.H.1.9]{BridsonHaefliger1999}, so \(\mathcal E\) is
word-hyperbolic.  The subgroup \(\Pi\) has finite index in
\(\mathcal E\), and hence is
word-hyperbolic as well.
\end{proof}

\subsection{Poincar\'e duality and outer automorphisms}

We use the closed aspherical topological four-manifold model of
\(\Pi\) to exclude splittings over virtually cyclic groups.  The
argument detects a splitting by a finitely supported path cocycle;
duality and the homological form of Shapiro's lemma then force its
cohomology group to vanish.

\begin{lemma}
\label{hyp:no-vc-splitting}
The group \(\Pi\) admits no nontrivial splitting over a virtually
cyclic subgroup.  It has one end.
\end{lemma}

\begin{proof}
The finite aspherical complex \(Q_2\) is a finite-dimensional
\(K(\Pi,1)\).  Its lifted cellular chain complex gives
\(\operatorname{cd}_{\mathbb Z}\Pi<\infty\).  A group of finite
integral cohomological dimension is torsion-free: if it contained a
subgroup of prime order, restricting a finite projective resolution
to that subgroup would contradict the nonzero integral cohomology of
a cyclic group of prime order in arbitrarily high degrees.  Thus
\(\Pi\) is torsion-free.  Also, \(\Pi\cap H\) has finite index in
the nonabelian free group \(H\), so it is a nonabelian free group.
In particular, \(\Pi\) is infinite and is not virtually cyclic.

Suppose that \(\Pi\) has a splitting as an amalgamated product or HNN
extension with virtually cyclic edge group \(U\).  Assume the
splitting is nontrivial, meaning that the
Bass--Serre tree \(\mathcal B\) has no global fixed vertex; the action
has no inversions.  Its unoriented edges form the \(\Pi\)-set
\(\Pi/U\).  Fix a vertex \(p\) of \(\mathcal B\), and define
\[
c(g)=\sum_{e\text{ on }[p,gp]}e
       \ \in\ \Ftwo[\Pi/U]\qquad(g\in\Pi),
\]
where the sum is empty when \(gp=p\).  All sums have finite support.
Concatenating the paths from \(p\) to \(gp\) and from \(gp\) to
\(ghp\), and canceling edges traversed twice, gives the path from
\(p\) to \(ghp\).  Over \(\Ftwo\) this says
\[
c(gh)=c(g)+g\,c(h).
\]
Thus \(c\) is a group-cohomology cocycle.

The orbit of \(p\) is unbounded.  Otherwise its bounded invariant
subtree would have a center fixed by \(\Pi\); the center is a vertex
or an edge midpoint, and in the latter case the absence of inversions
also fixes its endpoints.  Either case contradicts nontriviality.
Consequently
\[
|\supp c(g)|=d_{\mathcal B}(p,gp)
\]
is unbounded as \(g\) varies.  On the other hand, a coboundary
\(gq-q\), for \(q\in\Ftwo[\Pi/U]\), has support size at most
\(2|\supp q|\) for every \(g\).  Hence
\begin{equation}
[c]\ne0\quad\text{in}\quad H^1(\Pi;\Ftwo[\Pi/U]).
\label{hyp:path-class}
\end{equation}

Let \(M\) be the closed aspherical topological four-manifold with
fundamental group \(\Pi\).  Poincar\'e duality with local coefficients
on \(M\) applies to arbitrary \(\Ftwo\Pi\)-modules, including the
possibly infinite-dimensional permutation module \(\Ftwo[\Pi/U]\);
see \cite[Theorem~10.2]{Spanier1993}.  The orientation twist is trivial
over \(\Ftwo\).  Using the description of local coefficients by modules
over the fundamental group in \cite[Section~3.H]{Hatcher2002},
asphericity and duality therefore give
\[
H^1(\Pi;\Ftwo[\Pi/U])
 \cong H_3(\Pi;\Ftwo[\Pi/U]).
\]
The permutation module is the induced module
\(\Ftwo\Pi\otimes_{\Ftwo U}\Ftwo\), with trivial \(U\)-action on the
last factor.  Homological Shapiro's lemma
\cite[Lemma~6.3.2]{Weibel1994} gives
\begin{equation}
H_3(\Pi;\Ftwo[\Pi/U])\cong H_3(U;\Ftwo).
\label{hyp:shapiro}
\end{equation}
Indeed, if \(P_*\) is a free right \(\Ftwo\Pi\)-resolution of
\(\Ftwo\), then
\[
P_*\otimes_{\Ftwo\Pi}
   (\Ftwo\Pi\otimes_{\Ftwo U}\Ftwo)
 \cong P_*\otimes_{\Ftwo U}\Ftwo.
\]
The restriction of \(P_*\) is a free \(\Ftwo U\)-resolution, so the
right-hand complex computes the homology of \(U\).  This also makes
explicit why the induced, finite-support module is the one used here.

A torsion-free virtually cyclic group is either trivial or infinite
cyclic \cite[Sections~4.A.6.2 and~5.B.1]{Stallings1971}.  Thus \(U\) has a
classifying space that is a point or a circle, and
\(H_3(U;\Ftwo)=0\); see
\cite[Examples~6.1.3--6.1.4]{Weibel1994}.  Duality and
\eqref{hyp:shapiro} now contradict \eqref{hyp:path-class}.  This proves
nonsplitting.

Finally \(\Pi\) is finitely generated by
Proposition~\ref{hyp:word-hyperbolic} and is infinite.  Stallings'
ends theorem \cite[Section~1.B.6]{Stallings1971} says that a finitely
generated group with more than one end has a nontrivial splitting over
a finite subgroup.  Such a
subgroup is virtually cyclic, so the splitting just excluded would
exist if \(\Pi\) had more than one end.  Therefore \(\Pi\) has one
end.
\end{proof}

\begin{proposition}
\label{hyp:finite-out}
The outer automorphism group \(\Out(\Pi)\) is finite.
\end{proposition}

\begin{proof}
Bestvina and Feighn, using Paulin's construction of a limiting action
on an \(\mathbb R\)-tree \cite{Paulin1991}, prove that a word-hyperbolic
group with infinite outer automorphism group splits nontrivially over a virtually
cyclic subgroup \cite[Corollary~1.3]{BestvinaFeighn1995}.
The group \(\Pi\) is word-hyperbolic by
Proposition~\ref{hyp:word-hyperbolic}, while
Lemma~\ref{hyp:no-vc-splitting} excludes every such splitting.
The stated splitting theorem therefore implies that \(\Out(\Pi)\)
is finite.
\end{proof}
